\par\Needspace{4\baselineskip}
\begin{center}
\hypertarget{native-058-article040}{}\pdfbookmark[2]{40.}{article-040}40.
\end{center}
Für den dritten Schritt, die Bestimmung der wahren Anomalie, brauchen wir den Winkel $E$ selbst nicht. Es genügt $\operatorname{tang}\tfrac{1}{2}E$ oder genauer $\log\operatorname{tang}\tfrac{1}{2}E$. Dieser Schritt ließe sich bequem mit dem zweiten verbinden, wenn unsere Tafel unmittelbar den Logarithmus der Größe $\frac{\operatorname{tang}\tfrac{1}{2}E}{\BareRoot{A}}$ lieferte, die von $1$ um eine Größe zweiter Ordnung abweicht. Wir haben es jedoch vorgezogen, unsere Tafel etwas anders einzurichten. Bei geringerem Umfang lässt sich so dennoch viel bequemer interpolieren. Schreiben wir der Kürze halber $T$ für $\operatorname{tang}\tfrac{1}{2}E^2$, so lässt sich der in Art. 37 angegebene Wert von $A$, nämlich $\frac{15(E-\sin E)}{9E+\sin E}$, leicht umformen zu
\[
A=\frac{T-\tfrac{6}{5}T^2+\tfrac{9}{7}T^3-\tfrac{12}{9}T^4+\tfrac{15}{11}T^5-\text{usw.}}{1-\tfrac{6}{15}T+\tfrac{7}{25}T^2-\tfrac{8}{35}T^3+\tfrac{9}{45}T^4-\text{usw.}},
\]
\SourceLeaf{059}Das Bildungsgesetz der Reihe liegt auf der Hand. Durch Reihenreversion ergibt sich daraus
\[
\begin{aligned}
\frac{A}{T}&=1-\frac{4}{5}A+\frac{8}{175}A^2+\frac{8}{525}A^3\\
&\quad+\frac{1896}{336875}A^4+\frac{28744}{13138125}A^5+\text{usw.}\mbox{\HistoricalDisplayNoteMark}
\end{aligned}
\]
\noindent\HistoricalDisplayNoteText{Handschriftliche Bemerkung:]
\[
\begin{aligned}
\log\operatorname{hyp}\frac{A}{T}&=-\frac{4}{5}A-\frac{48}{175}A^2-\frac{104}{875}A^3-\cdot\cdot\cdot\\
&=-\frac{4}{5}T+\frac{64}{175}T^2-\frac{8}{35}T^3+\cdot\cdot\cdot
\end{aligned}
\]
\[
\log\frac{A}{T}=-2Am=-2Tm'
\]
\begin{center}
\begin{tabular}{c|c|c@{\hspace{2em}}c|c}
$A$ & $\log m$ & $\log n$ & $T$ & $\log m'$ \\
$0.000$ & $9.23984$ & $9.63778$ & $0.000$ & $9.23984$ \\
$0.010$ & $9.24134$ & $9.63909$ & $0.010$ & $9.23787$ \\
$0.020$ & $9.24284$ & $9.64041$ & $0.020$ & $9.23591$ \\
$0.030$ & $9.24435$ & $9.64172$ & &
\end{tabular}
\end{center}
\[
\begin{gathered}
\operatorname{tang}\tfrac{1}{2}v=M\gamma\operatorname{tang}\tfrac{1}{2}w,\qquad r=\frac{q\sec\tfrac{1}{2}v^2}{N},\\
\log M=Am,\qquad \log N=An.
\end{gathered}
\]
}
Setzen wir also $\frac{A}{T}=1-\frac{4}{5}A+C$, so ist $C$ eine Größe vierter Ordnung. Nehmen wir sie in unsere Tafel auf, so können wir unmittelbar von $A$ zu $v$ übergehen, und zwar nach der Formel
\[
\operatorname{tang}\tfrac{1}{2}v=\BareRoot{\frac{1+e}{1-e}}\,\mathord{.}\,\BareRoot{\frac{A}{1-\frac{4}{5}A+C}}=\frac{\gamma\operatorname{tang}\tfrac{1}{2}w}{\BareRoot{(1-\frac{4}{5}A+C)}},
\]
Dabei bezeichnet $\gamma$ die Konstante $\BareRoot{\frac{5+5e}{1+9e}}\cdot$ Zugleich erhalten wir so ein sehr bequemes Verfahren zur Berechnung des Radiusvektors. Es gilt nämlich (Art. 8, VI)
\[
\begin{aligned}
r&=\frac{q\cos\tfrac{1}{2}E^2}{\cos\tfrac{1}{2}v^2}=\frac{q}{(1+T)\cos\tfrac{1}{2}v^2}\\
&=\frac{(1-\frac{4}{5}A+C)q}{(1+\frac{1}{5}A+C)\cos\tfrac{1}{2}v^2}\cdot
\end{aligned}
\]
